\section{The automaton: states and rules}
\label{sec:automaton}

We define the automaton on the cell graph $G$ with the roles and labels of Section~\ref{sec:setting}, give its
complete state set and transition tables, and prove its motion and switches correct.

\subsection{The state of a cell}

A cell carries a static part, fixed by the initial configuration, and a dynamic part, updated each step. The
static part is the role (empty, track, or switch), the track-link labels, and, for a switch, its kind. The
dynamic part is one of four symbols and, for a switch, a one-bit selection. Table~\ref{tab:states} lists them.

\begin{table}[ht]
\centering
\begin{tabular}{@{}clp{8.2cm}@{}}
\toprule
symbol & name & meaning \\
\midrule
$\varnothing$ & empty & not part of the circuit, the quiescent background, stays empty \\
$C$ & clear track & a track or switch cell with no part of the locomotive on it \\
$H$ & head & the leading cell of the locomotive \\
$A$ & trailing & the cell just behind the head, the locomotive's tail \\
\midrule
$1,2$ & selection & a switch's selected branch, $1$ for branch $a$, $2$ for branch $b$ (switch cells only) \\
\bottomrule
\end{tabular}
\caption{The dynamic alphabet. Four cell symbols and, on a switch cell, one selection bit. The locomotive is
one cell in state $H$ with an adjacent cell in state $A$.}
\label{tab:states}
\end{table}

\subsection{The motion rules}

All cells update at once. On plain track the locomotive is governed by three rules, shown in
Table~\ref{tab:motion}. A head becomes the trailing cell, the trailing cell clears, and a clear track cell that
the head points into becomes the new head.

\begin{table}[ht]
\centering
\begin{tabular}{@{}cllc@{}}
\toprule
rule & a cell in state & with & becomes \\
\midrule
(M1) & $H$ & always & $A$ \\
(M2) & $C$ & the head points into it (the forward cell, Section~\ref{sec:switchrules}) & $H$ \\
(M3) & $A$ & always & $C$ \\
\bottomrule
\end{tabular}
\caption{The motion rules. Applied together each step, they advance the head-and-trailing pair forward by one
cell, and never backward, by Lemma~\ref{lem:direction}.}
\label{tab:motion}
\end{table}

\begin{lemma}[Direction]
\label{lem:direction}
On a stretch of plain track, where every track cell has exactly two track-neighbours, the only clear track cell
adjacent to the head is the cell ahead of the locomotive.
\end{lemma}

Let the head sit on a cell $c$, with track-neighbours $f$ ahead and $b$ behind. The cell behind is the trailing
cell, so $b$ is $A$, not $C$. The cell ahead, $f$, is $C$. So the unique $C$ track-neighbour of $c$ is $f$. \qed

\begin{lemma}[Motion]
\label{lem:motion}
If at time $t$ the locomotive is head $c$, trailing $b$, with $f$ ahead, then at time $t+1$ it is head $f$,
trailing $c$, and every other cell is unchanged.
\end{lemma}

By (M2) and Lemma~\ref{lem:direction}, $f$ becomes $H$. By (M1), $c$ becomes $A$. By (M3), $b$ becomes $C$. No
other cell is adjacent to the head, so nothing else changes. \qed

The argument uses only the graph and the roles, not the geometry, so the locomotive runs forward on any tiling.

\subsection{The switch rules}
\label{sec:switchrules}

A switch cell $s$ has trunk $u$, branches $a$ and $b$, kind, and selection. When the head occupies $s$, the
neighbour it came from holds the trailing cell $A$, so $s$ reads its entry side, and the forward cell and the
new selection are given by Table~\ref{tab:switch}.

\begin{table}[ht]
\centering
\begin{tabular}{@{}lllll@{}}
\toprule
kind & entered from & selection & forward cell & new selection \\
\midrule
any & branch $a$ or $b$ (passive) & $\sigma$ & trunk $u$ & $\sigma$ (memory sets it, below) \\
fixed & trunk (active) & $\sigma$ & $a$ if $\sigma{=}1$, else $b$ & $\sigma$ \\
flip-flop & trunk (active) & $1$ & $a$ & $2$ \\
flip-flop & trunk (active) & $2$ & $b$ & $1$ \\
memory & trunk (active) & $\sigma$ & $a$ if $\sigma{=}1$, else $b$ & $\sigma$ \\
\midrule
memory & a clear switch sees the head on branch $a$ & --- & --- & $1$ \\
memory & a clear switch sees the head on branch $b$ & --- & --- & $2$ \\
\bottomrule
\end{tabular}
\caption{The switch transition. The entry side is read from which track-neighbour holds the trailing cell. A
passive crossing goes to the trunk, an active crossing to the selected branch. The flip-flop flips on the
active crossing, the memory switch records the branch of a passive crossing one step before the head reaches
it.}
\label{tab:switch}
\end{table}

\begin{lemma}[Switches]
\label{lem:switches}
With Table~\ref{tab:switch}, a fixed switch always sends an actively entered locomotive to the same branch, a
flip-flop sends it to the selected branch and then selects the other, and a memory switch sends it to the
branch last entered passively. Every passive crossing goes to the trunk.
\end{lemma}

Each line of the table is the corresponding clause, and the entry side is unambiguous because, by
Lemma~\ref{lem:direction} applied at the switch with its three labelled track-neighbours, the trailing cell
sits on exactly the neighbour the head came from. \qed

\subsection{One fixed rule}

The rules of Tables~\ref{tab:motion} and~\ref{tab:switch} mention no tiling. They read a cell's role, its
dynamic symbol, the dynamic symbols of its neighbours, and a switch's labels and selection. The same rule runs
on every cell graph. By Lemmas~\ref{lem:direction} to~\ref{lem:switches} it realizes a track, a one-directional
locomotive, and the three switches, on any regular tiling. The one further rule, the track-building rule that
makes the registers self-extend, is the subject of the next section.
